% !TEX program = xelatex
%# -*- coding:utf-8 -*-
\documentclass[12pt,aspectratio=169,mathserif]{beamer}
% 设置为 Beamer 文档类型，设置字体为 12pt，长宽比为16:9，数学字体为 serif 风格

%%%%-----导入宏包-----%%%%
\usepackage{zju}            % 导入 zju 模板宏包
\usepackage{amsmath,amsfonts,amssymb,bm}
\usepackage{color}
\usepackage{graphicx,hyperref,url}
\usepackage{metalogo}
\usepackage{fontspec}
\usepackage{xeCJK}
\usepackage{tcolorbox}
\usepackage{tikz}
\usetikzlibrary{arrows.meta,positioning,calc,fit}
\beamertemplateballitem

%%%%------------------------%%%%%
\catcode`\。=\active
\newcommand{。}{．}
%%%%%%%%%%%%%%%%%%%%%

%%%%----首页信息设置----%%%%
\title[基于增量奇异值分解的偏微分方程模型降阶方法]{基于增量奇异值分解的偏微分方程模型降阶方法}
\author[Zhouyue Qian]{钱周越 \\ 指导老师：李雨文}
\institute[ZJU]{数学科学学院 \\ 浙江大学}
\date[Mar. 6th, 2026]{2026年3月6日}

%%%%----自定义命令----%%%%
\newcommand{\R}{\mathbb{R}}
\newcommand{\norm}[1]{\left\lVert #1\right\rVert}
\newcommand{\ipW}[2]{\left\langle #1, #2\right\rangle_{W}}

\begin{document}

%========================================================
\begin{frame}
  \titlepage
\end{frame}

%========================================================
\section{目录}
\begin{frame}
  \frametitle{目录}
  \tableofcontents
\end{frame}

%========================================================
\section{问题背景}

\begin{frame}
  \centering
  \Huge \insertsection
\end{frame}

%--------------------------------------------------------
\begin{frame}
\frametitle{参数化椭圆方程：many-query场景}

考虑含参数的一族椭圆型方程：
\[
-\nabla\cdot\big(a(x;\mu)\nabla u(x;\mu)\big)=f(x),\quad x\in\Omega,\ \mu\in\mathcal{P}\subset\R^p .
\]


有限元离散后得到大规模线性系统：
\[
A(\mu)u_h(\mu)=f(\mu),\quad A(\mu)\in\R^{N\times N},\ u_h(\mu)\in\R^{N},\ N\gg 1.
\]

\begin{itemize}
  \item 参数 $\mu$ 改变 $\Rightarrow$ 需反复求解不同的 $A(\mu)u=f$
  \item 难点：\textbf{高精度 FEM 解昂贵，重复求解成本不可接受}
\end{itemize}

\vspace{0.3em}
\textbf{目标：在保证精度的前提下显著降低多参数重复求解开销。}
\end{frame}


%--------------------------------------------------------
\begin{frame}
\frametitle{缩减基方法（RB）：Offline/Online 分解}

核心思想：用低维子空间逼近高维有限元解
\[
u_h(\mu)\approx V_r a(\mu),\quad V_r\in\R^{N\times r},\ r\ll N .
\]

Galerkin 投影到 $V_r$：
\[
V_r^T A(\mu)V_r\, a(\mu)=V_r^T f(\mu),
\]
在线阶段仅需解 $r\times r$ 系统。

\begin{itemize}
  \item \textbf{Offline：}构造缩减基 $V_r$（主要计算瓶颈）
  \item \textbf{Online：}快速求解并恢复 $u_r(\mu)=V_r a(\mu)$
\end{itemize}

\end{frame}

%--------------------------------------------------------
\begin{frame}
\frametitle{POD/SVD：最优低秩子空间与“快照矩阵”}

收集 $m$ 个参数样本对应的 FEM 快照：
\[
U=[u_h(\mu_1),\dots,u_h(\mu_m)]\in\R^{N\times m}.
\]

做 SVD（或加权 core SVD）得到 POD 基：
\[
U=\Phi\Sigma\Psi^T,\quad V_r:=\Phi_{(:,1:r)}.
\]

截断误差（Frobenius 范数意义下的最优性）：
\[
\|U-U_r\|_F^2=\sum_{i>r}\sigma_i^2,\quad U_r=\Phi_r\Sigma_r\Psi_r^T.
\]

\end{frame}

%========================================================
\section{问题分析}

\begin{frame}
  \centering
  \Huge \insertsection
\end{frame}

%--------------------------------------------------------
\begin{frame}
\frametitle{为什么 RB 要与增量 SVD 结合？}

RB 离线阶段经常呈\textbf{迭代/自适应}结构：
\begin{itemize}
  \item 参数点逐步采样（随机/拉丁超立方/贪婪）
  \item 快照矩阵 $U$ \textbf{按列增长}：$U_{\ell+1}=[U_\ell\ \ u_{\ell+1}]$
  \item 基空间 $V_r$ 需要持续更新
\end{itemize}

若每次新增快照都重新 SVD：
\[
\text{batch SVD 需要一次性存储/处理大矩阵，重复分解代价高。}
\]

\end{frame}

%--------------------------------------------------------
\begin{frame}
\frametitle{增量 SVD：以加权内积为例}

在 FEM 场景中，自然内积常为加权形式（质量矩阵 $W=M$）：
\[
\ipW{u}{v}=u^T W v.
\]

已知 $U_\ell$ 的（截断）core SVD：
\[
U_\ell \approx Q\Sigma R^T,\quad Q^T W Q=I_k,\ R^TR=I_k.
\]

新快照 $u_{\ell+1}$ 的投影与残差：
\[
d=Q^T(Wu_{\ell+1}),\quad e=u_{\ell+1}-Qd,\quad p=\sqrt{e^TWe}.
\]

构造仅 $(k+1)\times(k+1)$ 的小矩阵：
\[
Y=\begin{bmatrix}\Sigma & d\\ 0 & p\end{bmatrix},
\]
对 $Y$ 做 SVD 即可更新主子空间，并按阈值截断控制秩增长。

\end{frame}

%--------------------------------------------------------
\begin{frame}
\frametitle{增量 SVD：子空间更新与截断}

若 $p>0$，定义归一化方向
\[
\hat e = \frac{e}{p}.
\]

对小矩阵做 SVD：
\[
Y=\tilde U \tilde\Sigma \tilde V^T.
\]

则更新后的分解为
\[
Q^+=[Q\ \hat e]\tilde U,\quad   \Sigma^+=\tilde\Sigma,\quad  R^+=
\begin{bmatrix}
R & 0\\
0 & 1
\end{bmatrix}
\tilde V.
\]

\textbf{截断策略：}

\begin{itemize}
  \item 若 $\sigma_{k+1}<\texttt{tol}$，舍弃对应奇异方向
  \item 或若 $p<\texttt{tol}$，认为新快照已在当前子空间内
\end{itemize}

\end{frame}

%--------------------------------------------------------
\begin{frame}
\frametitle{增量 SVD 的优势}

\begin{itemize}
  \item \textbf{计算优势：}避免反复 batch SVD（随着 $\ell$ 增大尤其明显）
  \item \textbf{内存优势：}可流式处理，不必长期保存全部快照
  \item \textbf{结构可行性：}在典型 RB 场景中快照矩阵常呈现较快奇异值衰减，从而低秩近似有效
  \item \textbf{可控性：}截断策略（按 $\sigma_{k}$ 或残差 $p$）抑制无效维度增长
\end{itemize}
\end{frame}

%--------------------------------------------------------
\begin{frame}
\frametitle{增量SVD的数值稳定性}

传统增量SVD问题：

\begin{itemize}
\item 正交矩阵连乘导致正交性退化
\item 需频繁重正交
\item 影响奇异值精度
\end{itemize}



研究动机：

\begin{itemize}
 \item 是否可以设计一种新的增量 SVD 更新策略，在保证数值稳定性的同时， \textbf{避免多次频繁的重正交化操作}
\end{itemize}

\end{frame}

%========================================================
\section{研究方案}

\begin{frame}
  \centering
  \Huge \insertsection
\end{frame}

%--------------------------------------------------------
\begin{frame}
\frametitle{整体流程}

\centering
\begin{tikzpicture}[node distance=10mm, every node/.style={font=\small}, >=Latex]
  \node[draw, rounded corners, align=center, minimum width=28mm, minimum height=8mm] (pde)
    {参数化椭圆PDE\\ $A(\mu)u=f$};
  \node[draw, rounded corners, align=center, right=of pde, minimum width=28mm, minimum height=8mm] (fem)
    {高精度FEM\\求解快照 $u_h(\mu_i)$};
  \node[draw, rounded corners, align=center, right=of fem, minimum width=34mm, minimum height=8mm] (isvd)
  {逐步生成快照 $u_h(\mu_i)$\\ 流式增量 core SVD\\更新 $Q,\Sigma,R$\\（无需存全量 $U$）};
  \node[draw, rounded corners, align=center, below=of isvd, minimum width=34mm, minimum height=8mm] (basis)
    {提取POD/RB基\\ $V_r=Q(:,1:r)$};
  \node[draw, rounded corners, align=center, left=of basis, minimum width=34mm, minimum height=8mm] (online)
    {Online：投影求解\\ $V_r^TA(\mu)V_r a=V_r^Tf$};
  \node[draw, rounded corners, align=center, left=of online, minimum width=34mm, minimum height=8mm] (recover)
    {恢复近似解\\ $u_r(\mu)=V_r a(\mu)$};

  \draw[->] (pde) -- (fem);
  \draw[->] (fem) -- (isvd);
  \draw[->] (isvd) -- (basis);
  \draw[->] (basis) -- (online);
  \draw[->] (online) -- (recover);

  \node[above=2mm of fem, font=\scriptsize] {Offline：少量样本};
  \node[above=2mm of isvd, font=\scriptsize] {Offline：逐步更新};
  \node[below=2mm of online, font=\scriptsize] {Online：many-query};
\end{tikzpicture}

\end{frame}

%--------------------------------------------------------
\begin{frame}
\frametitle{研究内容}

\begin{enumerate}
  \item \textbf{模型与数据：}选取典型参数化椭圆模型（如变系数 Poisson），程序实现 FEM 生成快照矩阵

  \item \textbf{基构造：}实现 POD-RB（batch SVD）作为 baseline

  \item \textbf{增量算法：}在加权内积框架下实现增量 SVD（core SVD），从理论和算法层面探索新的增量 SVD 更新策略，以减少频繁重正交带来的计算开销

  \item \textbf{截断策略：}残差阈值 $p<\texttt{tol}$ 与奇异值阈值 $\sigma_k<\texttt{tol}$ 的秩控制

  \item \textbf{验证：}通过数值实验系统比较不同方法的性能，包括误差、正交性、计算速度和内存占用（batch SVD 与增量 SVD，对比不同增量策略）
\end{enumerate}

\end{frame}

%--------------------------------------------------------
\begin{frame}
\frametitle{研究设计：系统对比}

围绕“基构造阶段”进行系统对比：

\vspace{0.3cm}

\textbf{对比对象：}

\begin{itemize}
  \item \textbf{Baseline：}batch POD（全量 SVD）
  \item \textbf{方法1：}经典增量 SVD（含必要重正交）
  \item \textbf{方法2：}改进增量 SVD
\end{itemize}

\vspace{0.3cm}

\textbf{统一条件：}

\begin{itemize}
  \item 相同快照顺序与参数采样
  \item 相同加权内积 $W$
  \item 相同截断阈值（$\texttt{tol}$）
\end{itemize}

\vspace{0.3cm}

目标：分析不同算法在“数值稳定性与效率”上的差异。

\end{frame}

%--------------------------------------------------------

\begin{frame}
\frametitle{研究设计：系统验证}

\begin{columns}

%---------------- LEFT ----------------
\column{0.5\textwidth}

\textbf{1. 精度}

\[
\frac{\|u_h(\mu)-u_r(\mu)\|_W}{\|u_h(\mu)\|_W}
\]

比较不同基构造方式下的降阶误差。

\vspace{0.5cm}

\textbf{2. 正交性}

\[
\|I-Q^T W Q\|_F
\]

衡量增量更新过程中正交结构的保持程度。

%---------------- RIGHT ----------------
\column{0.5\textwidth}

\textbf{3. 离线效率}

\begin{itemize}
\item 基构造总时间
\item 快照逐步增长时的累计耗时
\end{itemize}

\vspace{0.5cm}

\textbf{4. 内存占用}

\begin{itemize}
\item 峰值内存
\item 是否需保存完整快照矩阵
\end{itemize}

\end{columns}

\end{frame}

%--------------------------------------------------------
\begin{frame}
\frametitle{时间规划（2025.11--2026.05）}

\textbf{阶段一：理论准备（2025.11--2025.12）}

\begin{itemize}
  \item 梳理 RB / POD 框架与误差度量
  \item 研读 增量 SVD 稳定性分析
\end{itemize}

\textbf{阶段二：程序实现（2026.01--2026.03）}

\begin{itemize}
  \item 实现变系数 Poisson 的 FEM 算例与快照生成
  \item 实现 batch POD-RB（baseline）以及加权内积下的增量 core SVD
  \item 构建正交性监测与截断策略
\end{itemize}

\textbf{阶段三：对比优化（2026.04）}

\begin{itemize}
  \item batch vs 增量；经典增量 vs 改进增量
  \item 精度、正交性、效率、内存系统评估
\end{itemize}

\textbf{阶段四：成果整理（2026.05）}

\begin{itemize}
  \item 总结实验结果，并完成论文与答辩材料
\end{itemize}

\end{frame}
%========================================================
\section{预期成效}

\begin{frame}
  \centering
  \Huge \insertsection
\end{frame}

%--------------------------------------------------------
\begin{frame}
\frametitle{预期成效}

\begin{itemize}

\item \textbf{精度：}在相同截断阶数 $r$ 下，达到与 batch POD-RB 相当的降阶精度

\item \textbf{离线效率：}通过增量 SVD 更新替代重复 batch SVD，减少频繁重正交带来的计算开销，从而降低基构造成本

\item \textbf{算法改进：}参考相关理论分析，探索新的增量 SVD 更新策略，在保持数值稳定性的同时减少正交矩阵的反复更新

\item \textbf{内存效率：}支持快照矩阵的流式处理，避免存储全部快照数据，适用于更大规模参数采样

\item \textbf{稳定性：}在 FEM 加权内积框架下维持良好的正交性，使降阶基构造与误差评估更加可靠

\end{itemize}

\end{frame}

%========================================================
\section{参考文献}

\begin{frame}{参考文献}
\footnotesize
\begin{thebibliography}{99}

\bibitem{zhang2022}
Y. Zhang,
An answer to an open question in the incremental SVD,
arXiv:2204.05398, 2022.

\bibitem{brand2006}
M. Brand,
Fast low-rank modifications of the thin singular value decomposition,
Linear Algebra Appl., 415 (2006), 20--30.

\bibitem{brand2002}
M. Brand,
Incremental SVD of uncertain data with missing values,
ECCV, 2002.

\bibitem{hesthaven2015}
J. S. Hesthaven, G. Rozza, B. Stamm,
Certified Reduced Basis Methods for Parametrized PDEs,
Springer, 2015.

\bibitem{patera_rozza}
A. T. Patera, G. Rozza,
Reduced Basis Approximation and A Posteriori Error Estimation for Parametrized PDEs.

\end{thebibliography}
\end{frame}

\end{document}